
A neural network encoder trained exclusively on small MLP and CNN checkpoints --- without ever processing a transformer --- can predict ViT model accuracy approximately three times better than a state-of-the-art weight-space self-supervised learning method. The mechanism is not a more expressive symmetry group. It is the choice of coordinate system.

This result emerges from the intersection of two research agendas: how to represent neural network weights usefully, and whether representations learned from one architecture family transfer to another. Model zoo analysis --- predicting properties of arbitrary trained networks without retraining --- is increasingly relevant for model selection, architecture search, and continual learning at scale. Existing methods tacitly assume that test models share the architecture family of training models. When that assumption fails, they fail in a specific, diagnosable way.

The SANE method [Schurholt2024Towards] --- which achieves $R^{2}=0.72$ in its intended within-architecture setting --- represents each neural network as a flattened sequence of weight scalars, agnostic to network topology. Applied to ViT checkpoints after training on MLP/CNN checkpoints, SANE maps all inputs to near-identical latent codes (latent $std\approx 0.0014)$. The representation has collapsed: every ViT model produces the same embedding regardless of its accuracy. This is not degradation in quality; it is a fundamental failure of the representation substrate under cross-architecture distribution shift.

The failure mode suggests the solution. Flat tokenization discards the relational structure of a neural network --- the computational graph of neurons connected by scalar weights. This structure is architecture-agnostic: a ViT attention projection layer, a CNN convolutional filter, and an MLP linear layer are all, at bottom, directed graphs of neurons and edges. Encoding weights in terms of graph structure rather than flat sequences provides a consistent vocabulary across architectures. A permutation-equivariant GNN trained on this representation processes ViT weight graphs even after training only on CNN weight graphs, because the graph schema is shared.

The central insight of this work is that \textbf{the directed computational graph schema (node=neuron, edge=weight) provides an architecture-agnostic coordinate system that enables SSL representations trained on MLP/CNN model zoos to transfer directly to ViT model zoo property prediction without ViT training data.} The specific symmetry group --- whether to additionally enforce scale equivariance --- is secondary. Counter to predictions from supervised same-architecture settings [Kalogeropoulos2024Scale], scale+permutation equivariance (EquiSSL) underperforms permutation-only equivariance (EquiSSL-perm) in the ViT cross-architecture SSL setting. We hypothesize that LayerNorm in ViT architectures fixes the scale degree of freedom that scale equivariance targets, making scale invariance an incorrect inductive bias for ViT weight encoding (Section 3.5). This hypothesis is empirically motivated; causal confirmation requires additional experiments on non-LayerNorm target architectures.

This paper makes four contributions:

\textbf{(1)} We identify and characterize the latent collapse failure mode of flat tokenization under cross-architecture SSL. On 53 real ViT-S/16 ImageNet checkpoints, SANE achieves $R^{2}=0.072$ with latent $std\approx 0.0014$ --- a mechanistic failure arising because flat tokenization provides no relational structure to distinguish ViT attention weights from MLP weights. SANE is not a generally poor method: its $R^{2}=0.072$ reflects the cross-architecture distribution shift failure specifically, not a deficiency in its intended domain.

\textbf{(2)} We demonstrate that graph-based SSL with permutation equivariance achieves meaningful cross-architecture transfer without target-domain training data. EquiSSL-perm achieves $R^{2}=0.231$ on 53 held-out ViT-S/16 models --- a ~220\% improvement over SANE. These results are from a single seed; full statistical validation with multiple seeds and the complete ViT zoo is required.

\textbf{(3)} We find that scale equivariance does not improve over permutation-only equivariance for ViT cross-architecture transfer, providing empirical support for the LayerNorm gauge-fixing hypothesis. This is a negative result with mechanistic interpretation, constraining when scale equivariance is a beneficial inductive bias.

\textbf{(4)} We diagnose why graph-based latent interpolation fails as a model editing primitive. Latent interpolation achieves acc=0.100 (exactly random chance for CIFAR-10), significantly worse than weight-space averaging (acc=0.125; t=-24.52, $p\approx 9\times 10^{-88}$, Cohen's d=-1.10). The failure is attributable to a decoder that reconstructs edge attribute statistics vectors rather than functional weight tensors --- a decoder design limitation, not a latent space quality issue.

